\section{The Duty to Give Reasons and the Procedural Regulation}

The enforcement model of the \gls{gdpr} presents a complex network of horizontal, vertical and diagonal relationships between the \gls{lsa}, the \glspl{sac}, and the \gls{edpb},\autocite[461]{hofmannmustertScholtenHandbook2023} which are all meant to collaborate in the enforcement of the \gls{gdpr} and ultimately contribute to the assessments included in the final decision of the \gls{lsa}.
As discussed in section 2.3.2, administrations have a duty to provide the parties affected by their decisions (i) a concise but
clear and relevant indication of the main issues of law and of fact that have led to the decision, (ii) in enough detail to understand the reasoning behind the decision.

The performance of this duty is a necessity in the integrated system of enforcement of the \gls{gdpr}, as \glspl{sa} are required, at multiple steps of the procedure, to exchange their views and legal assessments with each other, as well as with the \gls{edpb}.
This duty, according to Craig, is crucial for three reasons: it allows the parties to understand the reasons for a decision affecting them, it allows the judiciary to review the legality of the decision, and it serves as an internal check on the administration, as it requires the administration to lay out its reasoning based on the facts and the law.\autocite[369--370]{craigProcess2018}

This chapter will explore the duty of administrations to give reasons in the context of the new cooperation mechanism introduced by the Procedural Regulation. 
This study will focus on Craig's third function of reason giving and the ways the new steps introduced by the Procedural Regulation will increase the opportunities for the \glspl{sa} to give reasons to each other, potentially leading to more peer pressure and more carefully reasoned enforcement outcomes.

\subsection{Reason-giving in the Pre-Reform Cooperation Mechanism}

The cooperation mechanism of the \gls{gdpr} already provided for some opportunities for the \glspl{sa} to exchange their views.
While the \gls{gdpr} did not explicitly state what information was `relevant' for the purposes of the cooperation mechanism, this has been interpreted by the \gls{edpb} as including the factual elements and legal issues relating to a case.\autocite[46]{edpbGuidelines022022}
Notably, \glspl{sa} use voluntary mutual assistance to discuss and solve general legal questions relating to the application of the regulation, thus contributing to a uniform undertsanding of the \gls{gdpr} and helping each other in interpreting the legal framework.\autocite[4.1.2]{edpbContributionGDPREvaluation2023}

According to the \gls{edpb}'s interpretation of Article 60(3) \gls{gdpr}, the \gls{lsa} already had duty to consider the \glspl{sac}'s views all along the procedure, and not just at when a draft decision was issued.\autocite[124--125]{edpbGuidelines022022}
This obligation stems from the duty of \glspl{sa} to reach consensus on an investigation from the earliest stages, thus creating a mutual exchange of views and legal assessments througout the procedure, to reach a decision that is consensual.\autocite[125--128]{edpbGuidelines022022}

Naturally, another moment where the \glspl{sa} could already exchange and challenge each other's views is the draft decision stage, with an \gls{rro}.
\glspl{rro}, as seen in section 3.2.1, must be `reasoned'.
In particular, that element requires the \gls{sa} to provide `sound and substantiated reasoning for its objection' in a `coherent, clear, precise and detailed' way, in particular with respect to the legal arguments and factual elements it is basing its assessment on.\autocite[17--19]{edpbGuidelines092020}
The \gls{sa} must demonstrate the significance of the risks posed by the draft decision, and how the change would lead to a different outcome in the legal findings of the case, linking the consequences of the draft decision to the significance of the risks to the fundamental rights of data subjects.\autocite[16--19]{edpbGuidelines092020}

In case the objections are not followed, the \gls{edpb} will have to provide a binding decision on the matter,\autocite[65(1)(a)]{GDPR} which creates a further opportunity for the \glspl{sa} to exchange views and challenge each other's reasoning.
The decision of the \gls{edpb}, too, must be reasoned.\autocite[65(2)]{GDPR}
Similar requirements for reason-giving in this integrated system appear when requesting opinions, or triggering the urgency procedure.\autocite[66(1)--(3)]{GDPR}

Looking at Craig's third function of reason-giving, these moments may function to produce more carefully reasoned decisions, as the \glspl{sa} are required to confront with their peers.
This mechanism fails however when the informational asymetries between the \gls{lsa} and the \glspl{sac} do not allow the \glspl{sac} to meaningfully challenge the \gls{lsa}'s reasoning and form their own opinion,\autocite[504]{mustertProceduralisingEnforcement2026} as seen in section 3.2 and 4.1.

\subsection{Reason-Giving in the Post-Reform Mechanism}
\subsubsection{The Cooperation File and the Ability of the SAs to Formulate their Views}

The first step that may improve the ability of the \glspl{sac} to formulate their views is the introduction of the cooperation file.
The list provided by Article 9 provides explicitly for the inclusion of the `steps taken and legal analysis carried out aiming to determine whether there has been an infringement of [the \gls{gdpr}]'.\autocite[9(1)(i)]{ProceduralRegulation}
A centralized file containing more relevant information on a case may allow the \glspl{sac} to better understand the reasons of the \gls{lsa} and enable them to effectively steer enforcement in a different direction.

\subsubsection{New Procedural Steps Requiring Reason-Giving}

\paragraph{The Summary of Key Issues Stage}

A new crucial step introduced by the Procedural Regulation is the summary of key issues, which is meant to set the scope of the investigation by identifying the key legal issues and factual elements to be investigated.\autocite[10]{ProceduralRegulation}
This new procedural step requires an exchange of views between the \glspl{sa}, and provides a new opportunity to challenge the preliminary assessments made by the \gls{lsa} and to provide reasoned arguments for a different assessment.\autocite[10--11]{ProceduralRegulation}

Like the draft decision stage, the summary of key issues can be challenged by the \glspl{sac}, and when consensus cannot be reached, the \gls{lsa} must refer the matter to the \gls{edpb} for a binding decision.\autocite[11(6)]{ProceduralRegulation}
Dispute resolution, however, is only available at this stage in complaint-based cases, and not in own-volition inquiries, which limits the procedural opportunities to challenge the \gls{lsa}'s reasoning in those cases.\autocites[505]{mustertProceduralisingEnforcement2026}[11(6)]{ProceduralRegulation}

\paragraph{Other Procedural Steps Requiring Reason-Giving}

In addition to the summary of key issues, the Procedural Regulation makes other explicit references to \glspl{sa}'s duty to reason with each other along the procedure.

Article 23(3) specifically requires \glspl{rro} to be prepared in a manner that allows other \glspl{sa} to `prepare their positions' and enable the \gls{edpb} to resolve the dispute efficiently.\autocite[23(3)]{ProceduralRegulation}
Similarly, when referring matters to the \gls{edpb} for dispute resolution, the \gls{lsa} must provide the \glspl{rro} that it did not follow, as well as the reasons why it did not follow them.\autocite[27(3)]{ProceduralRegulation}
The \gls{edpb} must now decide whether the \glspl{rro} are relevant and reasoned and communicate it to the \glspl{sac} within 4 weeks of referral.\autocite[27(5)]{ProceduralRegulation}

The institution of these new procedural steps may provide the integrated system of enforcement more avenues to exchange views and test the logic of the legal assessments made.\autocite[1703]{liShoulderEnforcementEuropean2022}

\subsubsection{Early Resolution}

The early resolution stage presents an involution compared to the previous system of cooperation.
As seen in section 4.2, the early resolution procedure allows the \gls{lsa} to decide to close a case with a decision that cannot be challenged by a \gls{rro}.\autocite[23(1)(c)]{ProceduralRegulation}
The reforms introduced to the cooperation flow by the Procedural Regulation are mostly inapplicable to early resolution.\autocite[5(6)]{ProceduralRegulation}
Worse still, the early resolution procedure allows the \gls{sa} who received the complaint to declare it `devoid of purpose' without ever involving other \glspl{sa} in the assessment of the complaint.\autocite[5(1)(a)]{ProceduralRegulation}

\subsubsection{The Role of the EDPB as a Forum for Reason-Giving}

The increase in opportunities for the \gls{edpb} to issue binding decisions, such as on the scope of an investigation or on the need to act in order to protect the rights of data subjects following a failure to act by the \gls{lsa}, also creates new moments where the \glspl{sac} can challenge the reasoning of the \gls{lsa} and provide their own legal assessments.\autocite[11(6), 12(5)]{ProceduralRegulation}

One advantage of resorting to procedures at the \gls{edpb} level more frequenty is the decisions of the \gls{edpb} are public.\autocite[39]{EDPBRulesofProcedure}
The use of dispute resolution may bring disagreements between \glspl{sa} to the surface earlier, such as with the summary of key issues.
This may provide more accountability to the public, exposing the reasoning of the \glspl{sa} to scrutiny and allowing for more pressure to be exerted on the system to deliver effective protection.\autocite[1702]{liShoulderEnforcementEuropean2022}

The \gls{edpb} however remains a coordinating body, and does not gain investigative powers of its own, except for being able to request information from the \glspl{sa}.\autocites[27(4)]{ProceduralRegulation}[25]{klozaIfItAint2026}
The \gls{edpb} will still only rule on the disagreements between the \glspl{sa}, without being able to address the asymmetries that led the \glspl{sac} to fail to meaningfully participate in the procedure.

\subsubsection{The Enforceability of the New Procedural Steps}

As seen in section 4.6, Article 14 of the Procedural Regulation states that non-compliance with the time limits imposed by the regulations does not affect the validity of the procedural steps or the final decision they lead to.\autocite[14]{ProceduralRegulation}
As Mustert noted, this risks making the new procedural tools function as a mere organizational benchmark rather than an enforceable procedural guarantee.\autocite[507]{mustertProceduralisingEnforcement2026}

The lack of enforceability of new procedural steps may reproduce the same issues of informational asymmetry present in the pre-reform system, hindering the ability of the \glspl{sac} to meaningfully participate to the decision-making process.

\subsection{The Capacity of DPAs to Formulate their Views}

As seen in section 4.5, \glspl{dpa} are often under-resourced and under-staffed,\autocite[4.4]{edpbContributionGDPREvaluation2023} which pushes them to create diverging prioritization strategies at the national level,\autocite[501]{mustertProceduralisingEnforcement2026} as well as using informal channels for cooperation to exchange information in order to more flexibly manage their increasing workload.\autocites[466--467]{hofmannmustertScholtenHandbook2023}[4.1.2]{edpbContributionGDPREvaluation2023}

The difficulties in the ability of \glspl{sa} to participate actively to cases do not only stem from a lack of procedural moments that require them to exchange views.
As seen in 5.1, there was already a duty to exchange views all along the cooperation procedure that could be inferred from the law, as interpreted by the \gls{edpb}.

If the structural issues of under-resourcing are not addressed, \gls{sa} may still resort to using more informal channels of cooperation, which do not carry the same legal obligations and time limits created by the Procedural Regulation.\autocite[467]{hofmannmustertScholtenHandbook2023}
\glspl{dpa} experience informal tools positively, showing a willingness to cooperate in the enforcement of the \gls{gdpr} uniformly across the \gls{eu}.\autocite[4.1.2]{edpbContributionGDPREvaluation2023}
Those informal tools would escape the structure of the Procedural Regulation, thus rendering the new procedural steps less effective.

\subsection{Interim Conclusion}

This chapter has explored whether the new steps introduced by the Procedural Regulation will increase the ability of the integrated system of enforcement of the \gls{gdpr} to produce more carefully reasoned decisions through the lens of Craig's third function of reason-giving and peer pressure.
It found that while the new procedural steps create more opportunities for \glspl{sa} to exchange views and challenge each other's reasoning, the lack of enforceability of the procedural steps, as well as the unaddressed structural issues of capacity for \glspl{sa} may still hinder the ability of \glspl{sac} to meaningfully participate in the decision-making process, thus reducing the ability to produce carefully reasoned assessments and decisions.